% 专业学位博士研究生中期考核示例。
% 复制本文件后，先修改下列基本信息，再填写各报告部分。
\documentclass[dp]{xdupgtm}
% 学号
\newcommand\studentNo{1101110071}
% 论文题目，可使用两行
\newcommand\thesisTitleOne{专业学位博士论文题目}
\newcommand\thesisTitleTwo{}
% 姓名、学位类别、领域（方向）、指导教师和学院
\newcommand\authorName{张三}
\newcommand\major{电子信息}
\newcommand\submajor{新一代电子信息技术}
\newcommand\advisorName{李四}
\newcommand\school{电子工程学院}
% 报告日期和选题来源
\newcommand\submitDate{2026年10月8日}
\newcommand\subjectSource{国家重点研发计划}
% 需要电子签名时取消注释，并将图片放入 figures/sign/。
% 专业学位博士中期评价页包含校内和校外导师签名位置。
% \renewcommand\xnjsqm{\sign{lisi}}
% \renewcommand\xnjsqmrq{2026年10月8日}
% \renewcommand\xyjsqm{\sign{lisi}}
% \renewcommand\xyjsqmrq{2026年10月8日}
% \renewcommand\zzqm{\sign{lisi}}
% \renewcommand\cyqm{\sign{lisi}\emptysign\sign{lisi}}
% \renewcommand\cyqmrq{2026年10月8日}

\begin{document}
% 选题来源只需在 xzlyj 环境外修改 subjectSource。
\begin{xzlyj}\end{xzlyj}
\begin{yjmbnr}
% 在这里填写学位论文研究目标及研究内容；应与开题报告保持一致。
\end{yjmbnr}
\begin{ywcgz}
% 在这里填写目前已完成的理论、实验研究内容及获得的结论。
\end{ywcgz}
\begin{bxfsm}
% 在这里填写与开题报告不相符的情况；没有变化时填写“无”。
\end{bxfsm}
\begin{gzhjh}
% 在这里填写下一步工作计划、拟完成的研究内容和需要解决的关键技术。
\end{gzhjh}
\begin{kycg}
% 在这里填写已发表、已录用或拟发表的论文及其他研究成果。
\end{kycg}
\begin{xnjsyj}
% 在这里填写校内指导教师评价意见。
\end{xnjsyj}
\begin{xyjsyj}
% 在这里填写校外指导教师评价意见。
\end{xyjsyj}
\begin{pyjl}
% 在这里填写中期总结报告评语及“通过”或“不通过”结论。
\end{pyjl}
\end{document}
